\worksheettitle
  {Чтение и запись многозначных чисел}
  {\modulemark{M04} Версия учителя \textbullet\ маршрут 25–40 минут}

\begin{teacherbox}[Паспорт]
  \textbf{Результат:} ученик читает число и записывает его цифрами, не теряя
  нулевые и неполные классы. \textbf{Предварительно:} M03.

  Лист намеренно содержит запас практики. База и контрольная точка образуют
  минимальный маршрут; закрепление и разбор ошибок можно выбирать по темпу.
\end{teacherbox}

\begin{checkpointbox}[Вход]
  (50804007=50\,|\,804\,|\,007). Группа класса тысяч (804), число единиц
  этого класса (804). Если ученик группирует слева или пишет последний класс
  как (7), сначала дайте M03-R.
\end{checkpointbox}

\section{Открытие чтения}

(214\,036\,508): двести четырнадцать миллионов тридцать шесть тысяч
пятьсот восемь. В группе (036) читается число 36, но позиции группы
сохраняются в записи.

Формы: 21 тысяча, 32 тысячи, 15 тысяч, 104 тысячи. Не превращайте этот
фрагмент в отдельный урок русского языка: задача --- корректно прочитать
готовую числовую запись.

\begin{teacherbox}[Ответы базового чтения]
  \begin{enumerate}
    \item шесть тысяч восемнадцать;
    \item сорок две тысячи триста пять;
    \item семьсот тысяч сорок;
    \item девять миллионов шесть тысяч двенадцать;
    \item триста пять миллионов восемьдесят.
  \end{enumerate}
\end{teacherbox}

\section{Открытие записи}

Главная опора --- таблица классов. Первый слева класс может быть неполным;
все последующие группы содержат ровно три позиции. Неназванный промежуточный
класс получает группу (000).

\begin{teacherbox}[Ответы базовой записи]
  \begin{enumerate}
    \item (28\,005);
    \item (603\,042);
    \item (7\,900\,008);
    \item (50\,000\,024);
    \item (201\,003\,070).
  \end{enumerate}
  Ошибка (8\,5\,42): группы тысяч и единиц не дополнены слева нулями.
  Верно (8\,|\,005\,|\,042=8\,005\,042).
\end{teacherbox}

\section{Как дозировать лист}

\begin{resultbox}[Три возможных маршрута]
  \begin{itemize}
    \item \textbf{Быстрый класс:} вход, открытие, по 2 задания базы, разбор
      ошибок, контрольная точка, M04\(\star\).
    \item \textbf{Средний темп:} вся база, 2–3 задания каждого блока
      закрепления, контрольная точка.
    \item \textbf{Нужна тренировка:} база с таблицей, выбранные задания
      закрепления, затем контроль. При потере классов --- M04-R.
  \end{itemize}
\end{resultbox}

\section{Закрепление: ответы}

Чтение:
\begin{enumerate}
  \item восемнадцать тысяч четыре;
  \item девяносто тысяч девяносто;
  \item четыреста семь тысяч один;
  \item двенадцать миллионов тридцать тысяч четыре;
  \item восемьдесят миллионов пять тысяч шестьсот;
  \item девятьсот шесть миллионов пятнадцать.
\end{enumerate}

Запись: (105\,014); (900\,009); (2\,040\,100); (70\,006\,002);
(400\,000\,004); (302\,500\,050).

Парная проверка: «шестьдесят один миллион восемь тысяч тридцать»;
(107\,020\,006).

\section{Ошибки и диагностика}

\begin{teacherbox}[Ответы]
  \begin{enumerate}
    \item (5\,000\,012): потерян нулевой класс тысяч.
    \item «Тридцать четыре миллиона восемьдесят»: нулевой класс тысяч не
      читают, (080) относится к единицам.
    \item (200\,007): класс единиц должен занимать три позиции.
    \item «Семь миллионов четырнадцать тысяч пять»: группа (005) содержит
      число 5, а не 500.
  \end{enumerate}
\end{teacherbox}

\section{Контрольная точка и развилка}

Ответы: «шестьдесят три миллиона четыреста семь»;
(90\,008\,015). Группы (008) и (015) занимают три позиции и сохраняют
разряды внутри классов.

\begin{resultbox}[Решение]
  \begin{itemize}
    \item \textbf{M05:} оба преобразования выполнены, нулевой класс при чтении
      пропущен, промежуточный класс и позиции при записи сохранены.
    \item \textbf{M04-R:} потерян класс, непервый класс записан одной-двумя
      цифрами или ученик читает группу (015) как 150.
    \item \textbf{M04\(\star\):} результат достигнут и объяснение опирается на
      классы и позиции.
  \end{itemize}
\end{resultbox}

После коррекции: прочитать (71\,000\,306), затем записать «сорок два
миллиона пять тысяч девять». Ожидается (42\,005\,009).
